\documentclass[../portafolio.tex]{subfiles}

% Solo agregue paquetes en el preámbulo de ../portafolio.tex

\begin{document}



%%%%%%%%%%%%%%%%%%%%%%%%%%%%%%%%%%%%%%%%%%%%%%%%%%%%%%%%%%%%%%%%%%%%%%%%%%%%%%%%
\section{Curve Fitting}   % ejemplo: Derivadas numéricas , introducción a git , 

%\hfill \textbf{Fecha de la actividad:} 04 de noviembre de 2022

\medskip

%---------------------------------------------------------------------------------

%-------------------------------------------------------------% Selección de evidencias
\subsection{Analysis of Voyager Mission Magnetic Field Data}

In this activity, magnetic field measurements from NASA's Voyager mission were analyzed using numerical fitting and interpolation techniques. The objective was to explore different approaches for modeling and reconstructing experimental data while comparing their respective characteristics and applications.

A linear model of the form
\[
p(x) = a_0 + a_1x,
\]
where $x = t - t_0$, was fitted to the dataset in order to identify the overall trend of the measurements. To perform this task, a custom polynomial fitting routine was implemented, as shown in Code~\ref{cod:lineal}. The fitting coefficients were obtained through a least-squares approach based on the normal equations.

\begin{listing}
\begin{pythoncode}
class polyfit(PolyClass.Poly):
    def __init__(self, x, y, grado=1):
        G = np.vander(x, grado + 1, increasing=True)
        
        M = G.T @ G
        z = G.T @ y
        
        self.coef = Resolver(M, z)
        self.chi2 = np.sum((G @ self.coef - y)**2)
\end{pythoncode}
\caption{Implementation of a polynomial least-squares fitting routine.}
\label{cod:lineal}
\end{listing}

The linear fit was then applied to the Voyager dataset using the code shown in Code~\ref{cod:ajustelineal}.

\begin{listing}
\begin{pythoncode}
x = xi - xi[0]
f = polyfit(x, yi)
\end{pythoncode}
\caption{Linear fit applied to the Voyager magnetic field data.}
\label{cod:ajustelineal}
\end{listing}

In addition to the linear regression, a cubic spline interpolation was performed. Unlike regression models, spline interpolation constructs piecewise cubic polynomials that pass through every data point while maintaining smoothness between intervals. The interpolation was implemented using the \texttt{cspline} function shown in Code~\ref{cod:interpolcsp}.

\begin{listing}
\begin{pythoncode}
g = cspline(xi, yi)
\end{pythoncode}
\caption{Cubic spline interpolation of the Voyager dataset.}
\label{cod:interpolcsp}
\end{listing}

The results of both approaches are presented in Figures~\ref{fig:ajustelineal} and~\ref{fig:cspeline}. Figure~\ref{fig:ajustelineal} shows the linear fit superimposed on the original measurements, while Figure~\ref{fig:cspeline} displays the smooth curve generated by the cubic spline interpolation.

\begin{figure}[h!]
\centering
\includegraphics[width=\textwidth]{ajustecurvas1.png}
\caption{Voyager magnetic field data with a linear least-squares fit.}
\label{fig:ajustelineal}
\end{figure}

\begin{figure}[h!]
\centering
\includegraphics[width=\textwidth]{ajustecurvas2.png}
\caption{Voyager magnetic field data reconstructed using cubic spline interpolation.}
\label{fig:cspeline}
\end{figure}

\subsection{Discussion}

This activity highlights the fundamental difference between regression and interpolation techniques. The linear fit captures the overall trend of the data by minimizing the fitting error, making it useful for identifying large-scale behavior and obtaining simplified models.

In contrast, cubic spline interpolation aims to reproduce the dataset exactly by constructing a smooth curve that passes through every observation. As a result, interpolation provides a more detailed representation of the measured values, whereas regression focuses on extracting underlying trends.

Through this exercise, practical experience was gained in scientific data analysis, least-squares optimization, interpolation methods, and the visualization of experimental datasets. These techniques are widely used in computational physics, data science, and quantitative research for the analysis and interpretation of observational data.





\end{document}
